%% kapa: exam=1 round=27 year=2016 period=1 form=A subject=real_estate subject_ko=부동산학원론 seq=8 number=88
%% source: Q-Net 감정평가사 기출문제 (cst00309) artlSeq=5169402 file=extracted/1008/5169402/제27회 감정평가사 1차 문제/제27회 감정평가사 1교시문제(A형).pdf
%% conversion: kordoc 4.24.51 -> tools/gen1.py (auto); review=proofread-sonnet; flags=none
\begin{question}{88}
부동산권리분석활동을 위한 자료의 조사ㆍ확인 및 분석에 관한 설명으로 옳은 것은?
\begin{choices}
  \item 공간정보의 구축 및 관리 등에 관한 법령상 지적도에 기재된 지목의 부호가 「공」으로 표기되어 있어, 분석대상부동산의 지목을 공장용지로 확인ㆍ판단하였다.
  \item 용수 또는 배수를 위하여 일정한 형태를 갖춘 인공적인 수로ㆍ둑 및 그 부속시설물의 부지와 자연의 유수가 있거나 있을 것으로 예상되는 소규모 수로부지를 공간정보의 구축 및 관리 등에 관한 법령상의 지목인 하천으로 확인ㆍ판단하였다.
  \item 부동산경매에서 경락허가결정 확정 후 경매대금을 완납한 때에 경락인은 등기를 하여야만 목적부동산의 소유권을 취득하는 것으로 확인ㆍ판단하였다.
  \item 국토의 계획 및 이용에 관한 법령상 토지이용계획확인서를 통해 건물의 소재지, 구조, 용도 등의 사실관계를 확인ㆍ판단하였다.
  \item 공간정보의 구축 및 관리 등에 관한 법령상 토지대장의 등록사항을 통해 토지소유자가 변경된 날과 그 원인을 확인ㆍ판단하였다.
\end{choices}
\end{question}
\answer{5}
